\chapter{Aspectos Conceituais}\label{cap:conceitos}
% ===========================================================================

\section{Oblivious Transfer (OT)}

\subsection{Definição}

Oblivious Transfer (OT) é um protocolo criptográfico entre duas partes:

\begin{itemize}
\item \emph{Sender} $S$, que possui um conjunto de mensagens
$X = \{x_0, x_1, \ldots, x_{n-1}\}$;
\item \emph{Receiver} $R$, que escolhe um índice $i \in \{0, \ldots, n-1\}$.
\end{itemize}

O protocolo garante que:

\begin{enumerate}
\item $R$ aprende apenas $x_i$;
\item $R$ não aprende nenhuma informação sobre $x_j$ para $j \neq i$;
\item $S$ não aprende qual índice $i$ foi escolhido.
\end{enumerate}

Formalmente, o objetivo é realizar:
\[
(x_0, \ldots, x_{n-1}),\ i \longrightarrow x_i
\]
sem violar as propriedades de privacidade.

\section{Variantes de OT}

\subsection{1-out-of-2 OT}

O caso mais simples ocorre quando o \emph{sender} possui duas mensagens:
\[
(x_0, x_1)
\]
O \emph{receiver} escolhe um bit:
\[
r \in \{0, 1\}
\]
Ao final do protocolo, o \emph{receiver} obtém:
\[
x_r
\]
sem aprender $x_{1-r}$, e o \emph{sender} não descobre o valor de $r$.

\subsection{1-out-of-n OT}

Generalização do caso anterior:

\begin{itemize}
\item \emph{Sender} possui $n$ mensagens;
\item \emph{Receiver} escolhe um índice $i \in \{0, \ldots, n-1\}$;
\end{itemize}

O \emph{receiver} obtém apenas $x_i$, mantendo as mesmas garantias de
privacidade.

\subsection{k-out-of-n OT}

Nesta variante:

\begin{itemize}
\item \emph{Sender} possui $n$ mensagens;
\item \emph{Receiver} escolhe um subconjunto $I \subseteq \{0, \ldots, n-1\}$
tal que $|I| = k$.
\end{itemize}

O \emph{receiver} aprende:
\[
\{x_i : i \in I\}
\]
sem obter informação sobre as demais mensagens.

\subsection{Generalized Oblivious Transfer (GOT)}

O Generalized Oblivious Transfer (GOT) é uma generalização do Oblivious
Transfer tradicional, introduzida por Ishai e Kushilevitz.
\cite{tassa2011}

\subsection{Definição}

Seja Alice portadora de um conjunto de mensagens
\[
M = \{M_1, M_2, \ldots, M_n\}.
\]

Seja $[n] = \{1, 2, \ldots, n\}$ o conjunto de índices.

Define-se uma coleção $\mathcal{A} \subseteq 2^{[n]}$ de subconjuntos
permitidos, chamada de coleção monotônica decrescente, isto é,
\[
S \in \mathcal{A} \text{ e } S' \subseteq S \implies S' \in \mathcal{A}.
\]

\subsection{Funcionalidade}

Bob escolhe um subconjunto $S \subseteq [n]$ tal que $S \in \mathcal{A}$. Ao
final do protocolo:

\begin{itemize}
\item Bob aprende exatamente as mensagens $\{M_i\}_{i \in S}$;
\item Bob não aprende nenhuma informação sobre $\{M_i\}_{i \notin S}$;
\item Alice não aprende nenhuma informação sobre o subconjunto $S$ escolhido
por Bob.
\end{itemize}

\subsection{Relação com OT Clássico}

\begin{itemize}
\item No 1-out-of-n OT:
\[ \mathcal{A} = \{S \subseteq [n] \mid |S| = 1\}. \]
\item No k-out-of-n OT:
\[ \mathcal{A} = \{S \subseteq [n] \mid |S| = k\}. \]
\item No GOT:
\[ \mathcal{A} \subseteq 2^{[n]} \]
pode ser qualquer coleção monotônica decrescente.
\end{itemize}

\subsection{Interpretação}

O GOT permite especificar políticas arbitrárias de acesso, desde que sejam
fechadas por inclusão (monotonicidade decrescente). Por exemplo:

\begin{itemize}
\item todos os subconjuntos de tamanho $\leq k$;
\item qualquer subconjunto que não contenha simultaneamente dois índices
específicos;
\item qualquer subconjunto contido em um conjunto autorizado $T$.
\end{itemize}

\subsection{Implementação}

Ishai e Kushilevitz mostraram que o GOT pode ser implementado a partir de
múltiplas execuções paralelas de protocolos 1-out-of-2 OT.

\section{OT Extension}

Executar muitos OTs base (baseados em criptografia de chave pública) é
computacionalmente custoso.

A técnica de OT Extension permite:

\begin{enumerate}
\item Executar um pequeno número de OTs base;
\item Expandir para milhares ou milhões de OTs usando apenas criptografia
simétrica.
\end{enumerate}

Isso reduz drasticamente o custo computacional, tornando prática a computação segura multipartidária (\emph{Secure Multi-Party Computation} -- MPC).

\section{Criptossistema de Rabin}

O criptossistema de Rabin é uma família de esquemas de criptografia de chave
pública baseada em uma função armadilha cuja segurança, assim como a do RSA (Rivest--Shamir--Adleman),
está relacionada à dificuldade da fatoração de inteiros.

A função armadilha de Rabin tem a vantagem de que inverter essa função foi
matematicamente provado ser tão difícil quanto fatorar inteiros, enquanto não
há prova semelhante conhecida para a função armadilha do RSA. Ela tem a
desvantagem de que cada saída da função de Rabin pode ser gerada por qualquer
uma de quatro entradas possíveis; se cada saída for um texto cifrado, é
necessária complexidade adicional na decriptação para identificar qual das
quatro entradas possíveis era o texto plano verdadeiro. Tentativas ingênuas de
contornar esse problema frequentemente possibilitam um ataque de texto cifrado
escolhido para recuperar a chave secreta ou, ao codificar redundância no
espaço do texto plano, invalidam a prova de segurança em relação à fatoração.

Esquemas de criptografia de chave pública baseados na função armadilha de
Rabin são usados principalmente como exemplos em livros-texto. Em contraste, o
RSA é a base de esquemas padrão de criptografia de chave pública, como o
RSAES-PKCS1-v1\_5 e o RSAES-OAEP, que são amplamente utilizados na prática.
\cite{rabin2005}

\subsection{Definição}

Sejam $p$ e $q$ primos grandes, escolhidos de forma que
\[
p \equiv q \equiv 3 \pmod{4},
\]
e:
\[
n = p \cdot q
\]
Uma mensagem $M$ é convertida em um número:
\[
m < n
\]
A encriptação é definida como:
\[
c = m^2 \bmod n
\]
O \emph{ciphertext} é $c$.

A mensagem $m$ pode ser recuperada a partir do texto cifrado $c$ calculando
sua raiz quadrada módulo $n$, como segue.

Calcule a raiz quadrada de $c$ módulo $p$ e módulo $q$ utilizando as
seguintes fórmulas:
\begin{align*}
m_p &= c^{\frac{1}{4}(p+1)} \bmod p \\
m_q &= c^{\frac{1}{4}(q+1)} \bmod q
\end{align*}
Essas fórmulas só são válidas porque $p \equiv q \equiv 3 \pmod 4$: nesse
caso $\frac{1}{4}(p+1)$ e $\frac{1}{4}(q+1)$ são inteiros e, pelo critério de
Euler, $m_p^2 \equiv c \pmod p$ e $m_q^2 \equiv c \pmod q$. Por isso a geração
de chaves deve impor essa condição.

Utilize o algoritmo de Euclides estendido para encontrar $y_p$ e $y_q$ tais
que:
\[
y_p \cdot p + y_q \cdot q = 1.
\]
Aplique o Teorema Chinês do Resto para encontrar as quatro raízes quadradas de
$c$ módulo $n$:
\begin{align*}
r_1 &= (y_p \cdot p \cdot m_q + y_q \cdot q \cdot m_p) \bmod n \\
r_2 &= n - r_1 \\
r_3 &= (y_p \cdot p \cdot m_q - y_q \cdot q \cdot m_p) \bmod n \\
r_4 &= n - r_3
\end{align*}

Um desses quatro valores é o texto plano original $m$, embora não seja
possível determinar qual deles é o correto sem informação adicional.
